\documentclass[11pt]{article}
\usepackage[margin=1in]{geometry}
\usepackage{amsmath,amssymb,amsthm}
\usepackage{physics}
\usepackage{enumitem}
\usepackage{hyperref}
\usepackage{titlesec}

\theoremstyle{definition}
\newtheorem{definition}{Definition}[section]
\newtheorem{example}{Example}[section]
\newtheorem{remark}{Remark}[section]
\newtheorem{proposition}{Proposition}[section]
\newtheorem{theorem}{Theorem}[section]

\titleformat{\section}{\large\bfseries}{Lecture 10.\thesection}{1em}{}

\title{\textbf{Lecture 10: Quantum Channels} \\ \large Understanding Quantum Information \& Computation}
\author{Based on the lectures by John Watrous (IBM Quantum)}
\date{}

\begin{document}
\maketitle
\tableofcontents
\newpage

%----------------------------------------------------------------------
\section{Overview}
%----------------------------------------------------------------------

In the simplified description of quantum information, the only allowable state changes are unitary operations.  In reality, quantum systems interact with their environments, leading to \textbf{noise} and \textbf{decoherence}.  This lecture introduces \textbf{quantum channels}---the most general physically realisable transformations on quantum states described by density matrices.

%----------------------------------------------------------------------
\section{Motivation}
%----------------------------------------------------------------------

A unitary operation $\rho\mapsto U\rho\,U^\dagger$ is deterministic and reversible.  Many physical processes are neither:
\begin{itemize}
    \item A qubit interacting with a thermal bath.
    \item Transmission of a photon through a lossy fibre.
    \item Discarding part of a composite system (partial trace).
\end{itemize}

All of these are described by quantum channels.

%----------------------------------------------------------------------
\section{Definition of a Quantum Channel}
%----------------------------------------------------------------------

\begin{definition}[Quantum channel]
A \textbf{quantum channel} is a linear map $\Phi$ from density matrices on an input system $X$ to density matrices on an output system $Y$ satisfying:
\begin{enumerate}
    \item \textbf{Completely positive (CP):} $\Phi\otimes\mathrm{id}_R$ maps positive semidefinite operators to positive semidefinite operators for every reference system $R$.
    \item \textbf{Trace-preserving (TP):} $\tr\bigl(\Phi(\rho)\bigr)=\tr(\rho)$ for all $\rho$.
\end{enumerate}
A map satisfying both conditions is called a \textbf{CPTP map}.
\end{definition}

\begin{remark}
Complete positivity is strictly stronger than positivity.  The transpose map $\rho\mapsto\rho^T$ is positive but \emph{not} completely positive---it fails when tensored with an identity on an entangled reference system.
\end{remark}

%----------------------------------------------------------------------
\section{Representations of Quantum Channels}
%----------------------------------------------------------------------

\subsection{Kraus (Operator-Sum) Representation}

\begin{theorem}
A linear map $\Phi$ is a quantum channel if and only if there exist operators $\{K_1,\dots,K_r\}$ (called \textbf{Kraus operators}) satisfying
\[
    \Phi(\rho) = \sum_{k=1}^{r}K_k\,\rho\,K_k^\dagger,
    \qquad\text{with}\quad
    \sum_{k=1}^{r}K_k^\dagger K_k = I.
\]
\end{theorem}

The trace-preserving condition $\sum_k K_k^\dagger K_k=I$ ensures $\tr(\Phi(\rho))=\tr(\rho)$.

\begin{example}[Unitary channel]
$\Phi(\rho)=U\rho\,U^\dagger$.  Single Kraus operator $K_1=U$.
\end{example}

\subsection{Stinespring (Unitary Dilation) Representation}

\begin{theorem}[Stinespring]
Every quantum channel $\Phi$ on system $X$ can be realised as:
\begin{enumerate}
    \item Introduce an environment system $E$ initialised in a fixed state $\ket{0}_E$.
    \item Apply a joint unitary $V$ on $XE$.
    \item Discard (trace out) the environment.
\end{enumerate}
That is,
\[
    \Phi(\rho) = \tr_E\!\bigl(V(\rho\otimes\ket{0}\!\bra{0}_E)V^\dagger\bigr).
\]
\end{theorem}

This is the physical picture: the system interacts unitarily with an environment, and the environment is then inaccessible.

\subsection{Choi Representation}

The \textbf{Choi matrix} of $\Phi$ is
\[
    J(\Phi) = \sum_{a,b}\Phi(\ket{a}\!\bra{b})\otimes\ket{a}\!\bra{b}.
\]
$\Phi$ is completely positive iff $J(\Phi)\succeq 0$, and trace-preserving iff $\tr_Y(J(\Phi))=I_X$.

%----------------------------------------------------------------------
\section{Important Examples of Quantum Channels}
%----------------------------------------------------------------------

\subsection{Depolarising Channel}

\[
    \Phi_p(\rho) = (1-p)\,\rho + \frac{p}{3}\bigl(X\rho X + Y\rho Y + Z\rho Z\bigr),
    \qquad p\in[0,1].
\]
Equivalently, $\Phi_p(\rho)=(1-\tfrac{4p}{3})\rho + \tfrac{4p}{3}\cdot\tfrac{I}{2}$.  With probability $p$, a random Pauli error occurs.

\subsection{Dephasing (Phase Damping) Channel}

\[
    \Phi_p(\rho) = (1-p)\,\rho + p\,Z\rho Z.
\]
Destroys off-diagonal coherences while leaving populations unchanged.

\subsection{Amplitude Damping Channel}

Models energy relaxation (e.g.\ spontaneous emission).  Kraus operators:
\[
    K_0 = \begin{pmatrix}1&0\\0&\sqrt{1-\gamma}\end{pmatrix},
    \qquad
    K_1 = \begin{pmatrix}0&\sqrt{\gamma}\\0&0\end{pmatrix},
    \qquad\gamma\in[0,1].
\]

\subsection{Erasure Channel}

With probability $p$, the qubit is replaced by a ``flag'' state $\ket{e}$ orthogonal to the qubit space, signalling that erasure occurred.

\subsection{Partial Trace as a Channel}

For a bipartite system $XY$, $\Phi(\rho_{XY})=\tr_Y(\rho_{XY})$ is a valid quantum channel from $XY$ to $X$.

%----------------------------------------------------------------------
\section{Properties of Quantum Channels}
%----------------------------------------------------------------------

\begin{itemize}
    \item \textbf{Composition:} the composition of two channels is a channel.
    \item \textbf{Convexity:} a probabilistic mixture $p\,\Phi_1+(1-p)\Phi_2$ of channels is a channel.
    \item \textbf{Tensor products:} $\Phi_1\otimes\Phi_2$ (independent channels on separate systems) is a channel.
    \item Channels map the set of density matrices into itself.
    \item Unitary channels are the only \emph{reversible} channels.
\end{itemize}

%----------------------------------------------------------------------
\section{Quantum Channels and the Simplified Description}
%----------------------------------------------------------------------

In the simplified description, the state-change rule is $\ket{\psi}\mapsto U\ket{\psi}$ (or $\rho\mapsto U\rho U^\dagger$).  This corresponds to a channel with a single Kraus operator.  The general description strictly extends this:

\begin{center}
\begin{tabular}{lcc}
\hline
& \textbf{Simplified} & \textbf{General} \\\hline
Operations & Unitary $U$ & Quantum channel $\Phi$ (CPTP) \\
Reversible? & Always & Only if unitary \\
Noise modelling & No & Yes \\
Discarding subsystems & Not described & Partial trace \\
\hline
\end{tabular}
\end{center}

%----------------------------------------------------------------------
\section{Summary}
%----------------------------------------------------------------------

\begin{itemize}
    \item A \textbf{quantum channel} (CPTP map) is the most general physically allowed transformation of quantum states.
    \item Equivalent representations: \textbf{Kraus operators}, \textbf{Stinespring dilation}, \textbf{Choi matrix}.
    \item Key examples: depolarising, dephasing, amplitude damping, erasure, partial trace.
    \item Unitary evolution is the special case of a channel with one Kraus operator; all other channels are irreversible and model noise.
\end{itemize}

\end{document}
